% Pro M. Aemilio Scauro --- headnote
% pro-scauro, 54 BC, Rome

\textit{For M. Aemilius Scaurus, delivered at Rome in
54 BC. The defendant was the son of the famous princeps
senatus of the same name; he had governed Sardinia as
propraetor and, on his return, was prosecuted de
repetundis --- for extortion in his province --- by Q.
Servilius Caepio and P. Valerius Triarius, even as he
was standing as a candidate for the consulship. The
case turned in large part on Sardinian witnesses, and
the prosecution made much of the death of a Sardinian
woman, the wife of one Aris, said to have killed
herself to escape Scaurus's advances. Cicero answers
the moral charges by reconstructing the squalid
domestic intrigue behind the woman's death, and answers
the witnesses themselves with notorious ethnic
invective: the Sardinians, he urges, are the disowned
offspring of Carthaginians and Africans, the
skin-clad ``pelliti'' whose word is worth nothing
against the dignity of a noble house. Scaurus was
defended by a team --- Cicero, Hortensius, and others
--- and was acquitted.}

\textit{The speech survives only in fragments. The
short lettered pieces (\textsection a--p, and again
\textsection 45a--46o) are isolated quotations
preserved by Asconius, who introduced the speech in his
commentary, and by the grammarians; some are no more
than a single clause. The continuous numbered stretch
(\textsection 2--51) comes from the Ambrosian
palimpsest, which preserves a long run of the argument
on the worthlessness of the prosecution's evidence: the
disputed methods of the inquiry (\textsection 23--30),
the political motive behind the trial in the rivalry of
Appius Claudius (\textsection 31--36), the assault on
the credibility of the Sardinian witnesses
(\textsection 38--44), and the closing appeal to the
defendant's father and to the Metelli, his maternal
line (\textsection 46o--50). The text breaks off
unfinished. Because of this fragmentary transmission
the section labels are not a continuous run of numbers
but the editorial sigla of the surviving pieces,
reproduced here in order.}
